\documentclass[12pt,titlepage,dvipdfmx,twoside,fleqn,a4paper]{jsarticle}
\usepackage{multicol}
\usepackage{amsmath,amssymb,amsthm,amsfonts}
\usepackage[driver=dvipdfm,top=20truemm,bottom=20truemm,left=20truemm,right=20truemm]{geometry}
\usepackage{fancyhdr}
\usepackage{ascmac}
\usepackage{mathtools}
\usepackage[dvipdfmx]{graphicx}
\usepackage{thmbox}
\usepackage{bm}
\usepackage{tikz}
\usetikzlibrary{cd,calc}
\usepackage{xcolor}
\usepackage{paralist}
\usepackage{titlesec}
\usepackage{tcolorbox}
\tcbuselibrary{raster,skins,breakable,theorems,xparse}
\usepackage[version=3]{mhchem}
\usepackage[dvipdfmx]{pict2e}
\usepackage{otf}
\usepackage[inline]{enumitem}
\usepackage{mleftright}
\usepackage{esint}
\usepackage{bigints}
\usepackage{varwidth}
\usepackage{ulem}
\usepackage{wrapfig}
\usepackage{pxrubrica}
\usepackage{url}
\usepackage{booktabs}
\usepackage{here}
\usepackage{fancyvrb}
\usepackage{xstring}

%CTANにない.sty
\usepackage{ascolorbox}



\makeatletter


\newtcbox{\card}{on line, colback=white, colframe=black, boxrule=0.6pt, arc=0pt, boxsep=0pt, left=2pt, right=2pt, top=3pt, bottom=3pt}

%その他の設定
\everymath{\displaystyle}

\newcommand{\kaitou}{【解答】}

\newcommand{\kan}{,\;}

% 数式モードでのみ , の後ろに少しスペースを入れる
\mathchardef\mathcommaorig=\mathcode`\,
\begingroup
  \catcode`\,=\active
  \gdef,{%
    \mathcommaorig
    \mskip4mu
  }%
\endgroup
\mathcode`\,="8000

\newcommand{\visible}{visible}

\newcommand{\youten}{\ascboxD*{要点}}

\newcommand{\kyoutyou}[1]{\uline{\textbf{#1}}}

\title{\empty}
\author{\empty}
\date{\empty}

% 日付の取得
% 1文字目から4文字目を \myyear に格納
\StrMid{\jobname}{1}{4}[\myyear]

% 6文字目から7文字目を \mymonth に格納
\StrMid{\jobname}{6}{7}[\mymonth]

% 9文字目から10文字目を \myday に格納
\StrMid{\jobname}{9}{10}[\myday]




%インデント等の長さ調整
\setlength{\parindent}{1em}
\setlength{\columnsep}{5mm}
\setlength{\mathindent}{23pt}

\setlist{nosep,topsep=0.3mm,itemsep=0.5mm,leftmargin=*,labelsep=0.3mm}

\thickmuskip=1.0\thickmuskip
\medmuskip=0.5\medmuskip
\thinmuskip=0.5\thinmuskip
\arraycolsep=0.5\arraycolsep
\AtBeginDocument{
  \abovedisplayskip     =0.5\abovedisplayskip
  \abovedisplayshortskip=0.5\abovedisplayshortskip
  \belowdisplayskip     =0.5\belowdisplayskip
  \belowdisplayshortskip=0.5\belowdisplayshortskip
}


%箇条書きの調整
\renewcommand{\theenumi}{\arabic{enumi}}
\renewcommand{\labelenumi}{(\theenumi)}
\renewcommand{\theenumii}{\roman{enumii}}
\renewcommand{\labelenumii}{(\theenumii)}
\renewcommand{\theenumiii}{\roman{enumii}-\alph{enumiii}}
\renewcommand{\labelenumiii}{(\theenumiii)}

\renewcommand{\labelitemi}{$\cdot$\;}



%分数の調整
\DeclareRobustCommand*\frac[2]{
   \mathchoice
      {\@frac\textstyle{3.5}{#1}{#2}}
      {{\displaystyle\@frac\textstyle{3.5}{#1}{#2}}}
      {\@frac\scriptstyle{1.5}{#1}{#2}}
      {\@frac\scriptscriptstyle{1.5}{#1}{#2}}}
\def\@frac#1#2#3#4{
    {\begingroup\mkern#2mu\lower.44ex\hbox{$#1#3$}\mkern#2mu\endgroup
    \@@over \mkern#2mu\lower-.1ex\hbox{$#1#4$}\mkern#2mu}}
\DeclareRobustCommand\dfrac[2]{\displaystyle\@frac\textstyle{3.5}{#1}{#2}}
\DeclareRobustCommand\tfrac[2]{\textstyle\@frac\textstyle{3.5}{#1}{#2}}



\makeatother



%セクション番号をローマ数字に
\renewcommand{\thesection}{\Roman{section}}
\renewcommand{\thesubsection}{\thesection-\arabic{subsection}}

%セクションの体裁
\titleformat{\section}{\normalsize\bfseries}{\thesection\quad}{0pt}{}
\titleformat{\subsection}{\normalsize\bfseries}{\thesubsection\quad}{0pt}{}
\titlespacing*{\section}{0pt}{1.5em}{0.5em}
\titlespacing*{\subsection}{0pt}{1em}{0.5em}

%ページスタイルの設定
\pagestyle{plain}



\begin{document}

\begin{center}
\vspace*{2em}
{\large Inverse Scattering filed theory}

\vspace{2em}

Kenjiro Kimura\\
Center for Mathematical and Data Sciences, Kobe University\\
Noriaki Kimura\\
Integral Geometry Science Inc.
\vspace{2em}
\end{center}

\section{Introduction}

The most important issue in the theoretical analysis for inverse wave scattering problem is that
it should be based on the realistic observation processes and methods. In the theory described
in this paper, both the wave transmitting point and the wave receiving point are arbitrarily set
on the curved surface to receive the wave scattering data, and the inside of the object is
visualized using the acquired data. Transmitting antennas and receiving ones are installed on a
curve in a two-dimensional cross section that moves in the x-axis direction. Assuming that the
number of transmitting antennas and the number of receiving antennas are both $N$, the time
waveforms of $N$ squared are recorded, and the image of measurement target should be
reconstructed using the data. In the overview of this theory, the scattering field function required
for visualization should satisfy the partial differential equation defined in 6-dimensional space
$(t,\, x,\, y_1,\, y_2,\, z_1,\, z_2)$ or 7-dimensional space \quad $(t,\, x_1,\, x_2,\, y_1,\, y_2,\, z_1,\, z_2)$ which corresponds to the
coordinates of transmitting points and receiving points extendedly-defined not only on the
surface but also the inside. The three-dimensional target image can be obtained by solving this
equation using transmitted data and received data on the curved surface as the boundary
condition and setting $t \to 0$,\; $x_1 \to x_2\,(=x)$,\; $z_2 \to z_1\,(=z)$,\; $y_1 \to y_2\,(=y)$.\;
This theoretical method is most suitable for computer analysis, and it is considered that it can be effectively
used for the research about non-destructive subsurface imaging.

\section{Inverse Scattering filed theory}
\subsection{Inverse scattering problem}

Under the condition that the wave frequency $\omega$ is constant, the wave transmitting point $\mathbf{r}_1$ and
the wave receiving point $\mathbf{r}_2$ freely move inside the cross-section $\mathrm{D}$ perpendicular to the $x$
direction with a certain constraint condition, when the data obtained at the above situation is
expressed for $G(\mathbf{r}_1,\, \mathbf{r}_2,\, \omega)$, this function should be related to the distribution of scattering points
in the region. FIG.1 shows a situation in which the wave transmitted from $\mathbf{r}_1$ is reflected at the
point $\boldsymbol{\xi}$ and returns to the receiving point $\mathbf{r}_2$. Since there are many scattering points in the region,
$G(\mathbf{r}_1,\, \mathbf{r}_2,\, \omega)$ can be expressed as follows.

%%% ページ2 %%%

\begin{equation}
G(\mathbf{r}_1,\, \mathbf{r}_2,\, \omega)
= \iiint_D \varphi\!\left(\mathbf{r}_1 \to \boldsymbol{\xi} \to \mathbf{r}_2,\, \omega\right) d\boldsymbol{\xi}
\tag{1-1}
\end{equation}

Here, $\varphi(\mathbf{r}_1 \to \boldsymbol{\xi} \to \mathbf{r}_2,\, \omega)$ indicates the signal of the
wave at the receiving point $\mathbf{r}_2$ when the wave is
transmitted from the $\mathbf{r}_1$ and scattered at point $\boldsymbol{\xi}$. The
constraint conditions imposed on the wave transmitting
point $\mathbf{r}_1$ and the wave receiving point $\mathbf{r}_2$ are as follows.
The constraint is that the $x$-coordinates of $\mathbf{r}_1$ and $\mathbf{r}_2$ are
always equal. This function $G(\mathbf{r}_1,\, \mathbf{r}_2,\, \omega)$ is used to
explain the theoretical structure of the inverse
scattering problem. Let $\mathrm{D}$ be a sub-region in the three-dimensional space, and let $\partial\mathrm{D}$ be its boundary. At
this situation, the function $G(\mathbf{r}_1,\, \mathbf{r}_2,\, \omega)$ is the solution
of the following differential equation inside the
domain $\mathrm{D}$.

\begin{figure}[H]
\centering
% ここにFigure 1の画像を挿入
% \includegraphics[width=0.5\textwidth]{fig1.png}
\vspace{6cm}
\caption{$\mathbf{r}_1$ and $\mathbf{r}_2$ move freely inside region $\mathrm{D}$, respectively. $G$ is the sum of the scattering signals from all points $\boldsymbol{\xi}$.}
\end{figure}

\begin{figure}[H]
\centering
% ここにFigure 2の画像を挿入
% \includegraphics[width=0.7\textwidth]{fig2.png}
\vspace{8cm}
\caption{Sensor array}
\end{figure}

%%% ページ3 %%%

\begin{equation}
L\!\left(\frac{\partial}{\partial t},\, \frac{\partial}{\partial\mathbf{r}_1},\, \frac{\partial}{\partial\mathbf{r}_2}\right) G(\mathbf{r}_1,\, \mathbf{r}_2,\, t) = 0
\tag{1-2}
\end{equation}

$\overline{G}(\mathbf{r}_1,\, \mathbf{r}_2,\, t)$ is a Fourier transform function of $G(\mathbf{r}_1,\, \mathbf{r}_2,\, \omega)$ with respect to $\omega$. Further, the value
of $G(\mathbf{r}_1,\, \mathbf{r}_2,\, \omega)$ at the boundary $\partial\mathrm{D}$ is a signal received by the antennas. The above equation
is solved under this boundary condition, and $\rho(\mathbf{r})$ is defined as follows from the result.

\begin{equation}
\rho(\mathbf{r}) = \lim_{t \to 0}\left[\mathrm{Tr}\!\left[\,\overline{G}(\mathbf{r}_1,\, \mathbf{r}_2,\, t)\,\right]\right]
= \overline{G}(\mathbf{r},\, \mathbf{r},\, 0)
\tag{1-3}
\end{equation}

This $\rho(\mathbf{r})$ is a function corresponding to the spatial distribution of the gradient of the
permittivity inside the region D using microwaves. In this scheme, deriving the differential
operator $L(\partial/\partial t,\, \partial/\partial\mathbf{r}_1,\, \partial/\partial\mathbf{r}_2)$ that appears here is the most esoteric problem.

\subsection{Derivation of $L\!\left(\partial/\partial t,\, \partial/\partial\mathbf{r}_1,\, \partial/\partial\mathbf{r}_2\right)$}

The method for deriving this differential operator is
be described as follows. A case where the sensor array
shown in FIG.\ 2 moves in the x-axis direction along a
curved surface while performing a rotational scan in the direction
perpendicular to the x-axis is considered. On an
arbitrary curve, the $y$ and $z$ coordinates of $\mathbf{r}_1$ and $\mathbf{r}_2$ are
not always equal.
Specifically, $\mathbf{r}_1 = (x,\, y_1,\, z_1)$, $\mathbf{r}_2 = (x,\, y_2,\, z_2)$. The function $G$ is defined as follows.

\begin{figure}[H]
\centering
% ここにFigure 3の画像を挿入
% \includegraphics[width=0.6\textwidth]{fig3.png}
\vspace{7cm}
\caption{Waves transmitted from $\mathrm{P}_1$ are reflected at point $\mathrm{P}$ and return to $\mathrm{P}_2$. $z_1$ and $z_2$ are optional. The measurement point moves on the cross-section curve $\mathrm{S}$.}
\end{figure}

\begin{equation}
G(\mathbf{r}_1,\, \mathbf{r}_2,\, \omega)
= \iiint_D \varphi\!\left(\mathbf{r}_1 \to \boldsymbol{\xi} \to \mathbf{r}_2,\, \omega\right) d\boldsymbol{\xi}
\tag{2-1}
\end{equation}

The equation that the function $G$ satisfies is derived. Here, $\omega = ck$. $c$ is the propagation velocity
and $k$ is the wave number, and the wavelength is $\lambda$, there is a relationship of $k = 2\pi/\lambda$. The

%%% ページ4 %%%

following function $\phi$ is introduced.

\begin{align}
\phi(x,\, y_1,\, y_2,\, z_1,\, z_2,\, \omega)
&= \iint_D \frac{e^{ik\rho_1}}{\rho_1}\,\frac{e^{ik\rho_2}}{\rho_2}\,\varepsilon(\xi,\, \eta,\, \zeta)\,d\xi\,d\eta\,d\zeta \notag\\
\rho_1 &= \sqrt{(x-\xi)^2 + (y_1-\eta)^2 + (z_1-\zeta)^2} \notag\\
\rho_2 &= \sqrt{(x-\xi)^2 + (y_2-\eta)^2 + (z_2-\zeta)^2}
\tag{2-2}
\end{align}

Here, the wave number is written as $k$, assuming that the time factor is proportional to
$\exp(-i\omega t)$. The function in the integrand in the above equation is $\varphi$ in equation (2-1).

\begin{equation}
\varphi = \frac{e^{ik\rho_1}}{\rho_1}\,\frac{e^{ik\rho_2}}{\rho_2}\,\varepsilon(\xi,\, \eta,\, \zeta)
\tag{2-3}
\end{equation}

A partial differential equation whose an asymptotic solution is equation (2-2) is derived. Here,
the higher-order terms of $1/\rho$ resulting from differentiations should be ignored. The
abbreviations for differentiations are defined as follows.

\begin{equation}
\frac{\partial}{\partial t} \to \partial_t,\quad
\frac{\partial}{\partial x} \to \partial_x,\quad
\frac{\partial}{\partial y_1} \to \partial_{y1},\quad
\frac{\partial}{\partial y_2} \to \partial_{y2},\quad
\frac{\partial}{\partial z_1} \to \partial_{z1},\quad
\frac{\partial}{\partial z_2} \to \partial_{z2}
\tag{2-4}
\end{equation}

As a result of some calculations, it is derived that $\phi$ satisfies the following equation.

\begin{equation}
\left[\frac{1}{4}\Delta_s^2 - (ik)^2\,\partial_x^2
- (\partial_{y1}^2 + \partial_{z1}^2)(\partial_{y2}^2 + \partial_{z2}^2)\right]\phi = 0
\tag{2-5}
\end{equation}

This equation is derived assuming a steady state, but it is not difficult to extend it to the unsteady
state. Therefore, the variables are replaced as follows.

\begin{equation}
-ik \to \frac{1}{c}\,\partial_t
\tag{2-6}
\end{equation}

Finally, the following equation is derived.

%%% ページ5 %%%

\begin{equation}
\left[\frac{1}{4}\Delta_s^2
- \frac{1}{c^2}\,\partial_t^2\,\partial_x^2
- (\partial_{y1}^2 + \partial_{z1}^2)(\partial_{y2}^2 + \partial_{z2}^2)\right]\phi = 0
\tag{2-7}
\end{equation}

Assuming that the time factor of $\phi$ is proportional to $\exp(-i\omega t)$, the solution of equation (2-7)
is derived. Firstly, multiple Fourier transforms are performed on $\phi$ for $t$, $x$, $y_1$, $y_2$.

\begin{equation}
\tilde{\phi}(k_x,\, k_{y_1},\, k_{y_2},\, z_1,\, z_2,\, \omega)
= \int_{-\infty}^{\infty} e^{i\omega t}\,dt
  \int_{-\infty}^{\infty} e^{ik_{y_1} y_1}\,dy_1
  \int_{-\infty}^{\infty} e^{ik_{y_2} y_2}\,dy_2
  \int_{-\infty}^{\infty} e^{ik_x x}\,\phi(x,\, y_1,\, y_2,\, z_1,\, z_2,\, t)\,dx
\tag{2-8}
\end{equation}

The following equations are obtained by describing the partial derivatives with respect to $z_1$
and $z_2$ as $D_{z1}$ and $D_{z2}$, respectively.

\begin{equation}
\{(D_{z1}^2 + D_{z2}^2 - k_x^2 - k_{y1}^2 - k_{y2}^2)^2
- 4k^2 k_x^2
- 4(D_{z1}^2 - k_{y1}^2)(D_{z2}^2 - k_{y2}^2)\}\,\tilde{\phi} = 0
\tag{2-9}
\end{equation}

When this equation is solved, there are two variables, $z_1$ and $z_2$. Therefore, the boundary value
problem does not hold unless the boundary condition is given in the region with one-dimensional degrees of freedom in the $(z_1,\, z_2)$ space with respect to the fixed coordinates
$(x,\, y_1,\, y_2)$, that is the coordinates $(k_x,\, k_{y1},\, k_{y2})$ in the wave number space. The boundary
condition acquired by the actual measurement is given only at one point $(f(y_1),\, f(y_2))$ in $(z_1,\, z_2)$ space. To resolve this important issue, the following fact should be taken into account. The
theory in this study, where $z_1$ is independent on $z_2$, should include a specific case where $z_1 = z$,
$z_2 = z$. The solution of equation (2-9) can be expressed as follows.

\begin{equation}
E(k_x,\, k_{y_1},\, k_{y_2},\, z_1,\, z_2)
= \exp(is_1 z_1)\exp(is_2 z_2)
\tag{2-10}
\end{equation}

This formula becomes as follows when $z_1 = z_2 = z$.

\begin{equation}
E(k_x,\, k_{y_1},\, k_{y_2},\, z_1,\, z_2)
= \exp\{i(s_1 + s_2)z\}
\tag{2-11}
\end{equation}

Substituting equation (2-10) into equation (2-9) gives the following equation.

%%% ページ6 %%%

\begin{equation}
(s_1^2 + s_2^2 + k_x^2 + k_{y1}^2 + k_{y2}^2)^2
- 4k^2 k_x^2
- 4(s_1^2 + k_{y1}^2)(s_2^2 + k_{y2}^2) = 0
\tag{2-12}
\end{equation}

Another equation is needed to derive each of $s_1$ and $s_2$. Here, from the above-mentioned
requirement for the consistency about the fact where this theory should include the specific case
$z_1 = z_2 = z$, $s_1$ and $s_2$ should satisfy the following relational expression.

\begin{equation}
s_1 + s_2
= \sqrt{\left(\sqrt{k^2 - k_{y1}^2} + \sqrt{k^2 - k_{y2}^2}\right)^2 - k_x^2}
\tag{2-13}
\end{equation}

From equations (2-12) and (2-13), $s_1\,(k_x,\, k_{y1},\, k_{y2})$ and $s_2\,(k_x,\, k_{y1},\, k_{y2})$ are derived as follows.

\begin{align}
s_1(k_x,\, k_{y1},\, k_{y2})
&= \frac{\sqrt{k^2 - k_{y1}^2}\;\sqrt{\left(\sqrt{k^2 - k_{y1}^2} + \sqrt{k^2 - k_{y2}^2}\right)^2 - k_x^2}}
{\sqrt{k^2 - k_{y1}^2} + \sqrt{k^2 - k_{y2}^2}}
\notag\\[6pt]
s_2(k_x,\, k_{y1},\, k_{y2})
&= \frac{\sqrt{k^2 - k_{y2}^2}\;\sqrt{\left(\sqrt{k^2 - k_{y1}^2} + \sqrt{k^2 - k_{y2}^2}\right)^2 - k_x^2}}
{\sqrt{k^2 - k_{y1}^2} + \sqrt{k^2 - k_{y2}^2}}
\tag{2-14}
\end{align}

Specific calculations are performed later, and the solution of equation (2-7) can be described as
follows using these $s_1\,(k_x,\, k_{y1},\, k_{y2})$ and $s_2\,(k_x,\, k_{y1},\, k_{y2})$.

\begin{align}
&\phi(x,\, y_1,\, y_2,\, z_1,\, z_2,\, k) \notag\\
&= \frac{1}{(2\pi)^3}
\int_{-\infty}^{\infty}\!\!\int_{-\infty}^{\infty}\!\!\int_{-\infty}^{\infty}
e^{-i(k_x x + k_{y_1} y_1 + k_{y_2} y_2)}\,
a(k_x,\, k_{y_1},\, k_{y_2})\,
e^{is_1(k_x,\,k_{y_1},\,k_{y_2})\,z_1}\,
e^{is_2(k_x,\,k_{y_1},\,k_{y_2})\,z_2}\,
dk_x\,dk_{y_1}\,dk_{y_2}
\tag{2-15}
\end{align}

The equation of the cross-section curve $\mathrm{S}$ with the $x$-coordinate fixed is assumed as follows.
Although not an essential assumption, for example, a paraboloid is assumed.

\begin{equation}
f(y) = \alpha y^2
\tag{2-16}
\end{equation}

The boundary conditions given on the cross-section curve $\mathrm{S}$ are as follows.

%%% ページ7 %%%

\begin{align}
&\phi(x,\, y_1,\, y_2,\, \alpha y_1^2,\, \alpha y_2^2,\, k) \notag\\
&= \frac{1}{(2\pi)^3}
\int_{-\infty}^{\infty}\!\!\int_{-\infty}^{\infty}\!\!\int_{-\infty}^{\infty}
e^{-i(k_x x + k_{y_1} y_1 + k_{y_2} y_2)}\,
a(k_x,\, k_{y_1},\, k_{y_2})\,
e^{i\alpha\{s_1(k_x,\,k_{y_1},\,k_{y_2})\,y_1^2 + s_2(k_x,\,k_{y_1},\,k_{y_2})\,y_2^2\}}\,
dk_x\,dk_{y_1}\,dk_{y_2}
\tag{2-17}
\end{align}

This equation is used to determine $a(k_x,\, k_{y1},\, k_{y2})$. The following abbreviated notation is adopted.

\begin{align}
a(\mathbf{k}) &= a(k_x,\, k_{y_1},\, k_{y_2}) \notag\\
s_1(\mathbf{k}) &= s_1(k_x,\, k_{y_1},\, k_{y_2}) \notag\\
s_2(\mathbf{k}) &= s_2(k_x,\, k_{y_1},\, k_{y_2})
\tag{2-18}
\end{align}

Equation (2-17) is the following integral equation for $a\,(k_x,\, k_{y1},\, k_{y2})$.

\begin{align}
&\phi(x,\, y_1,\, y_2,\, \alpha y_1^2,\, \alpha y_2^2,\, k) \notag\\
&= \frac{1}{(2\pi)^3}
\int_{-\infty}^{\infty}\!\!\int_{-\infty}^{\infty}\!\!\int_{-\infty}^{\infty}
e^{-i(k_x x + k_{y_1} y_1 + k_{y_2} y_2)}\,
a(\mathbf{k})\,
e^{i\alpha\{s_1(\mathbf{k})\,y_1^2 + s_2(\mathbf{k})\,y_2^2\}}\,
d\mathbf{k}
\tag{2-19}
\end{align}

%%% ページ8前半 %%%

When $a(\mathbf{k})$ can be obtained from this equation, visualization process become accomplished
from equation (2-15).

\begin{align}
&\phi(x,\, y_1,\, y_2,\, z_1,\, z_2,\, k) \notag\\
&= \frac{1}{(2\pi)^3}
\int_{-\infty}^{\infty}\!\!\int_{-\infty}^{\infty}\!\!\int_{-\infty}^{\infty}
e^{-i(k_x x + k_{y_1} y_1 + k_{y_2} y_2)}\,
a(\mathbf{k})\,
e^{is_1(\mathbf{k})\,z_1}\,
e^{is_2(\mathbf{k})\,z_2}\,
d\mathbf{k}
\tag{2-20}
\end{align}

In the above equation, $z_1 = z_2 = z$ is set, and then Fourier transform with respect to $k$ is
proceeded to obtain the following visualization function $\rho(\mathbf{r})$.

\begin{equation}
\rho(\mathbf{r}) = \lim_{t \to 0}\left[\frac{1}{2\pi}
\int_{-\infty}^{\infty} \phi(x,\, y,\, y,\, z,\, z,\, k)\,e^{-ickt}\,dk\right]
\tag{2-21}
\end{equation}

%%% ページ8後半 %%%

\subsection{Solution of integral equation}

$\Phi\,(k_x,\, y_I,\, y_J,\, k)$, which is Fourier transform of the data $\phi(x,\, y_I,\, y_J,\, z_I,\, z_J,\, t)$ measured at points $\mathrm{P}_I$
and $\mathrm{P}_J$ on the curved surface, is written as follows.

\begin{equation}
\Phi_{I,J}(k_x,\, y_I,\, y_J,\, k)
= \int_{-\infty}^{\infty}\!\!\int_{-\infty}^{\infty}
e^{-ikt - ik_x x}\,\phi(x,\, y_I,\, y_J,\, z_I,\, z_J,\, t)\,dt\,dx
\tag{3-1}
\end{equation}

\begin{figure}[H]
\centering
% ここにFigure 4の画像を挿入
% \includegraphics[width=0.7\textwidth]{fig4.png}
\vspace{8cm}
\caption{The wave transmitted from $\mathrm{P}_1$ is reflected at point $\mathrm{P}$ and returns to $\mathrm{P}_2$. The boundary surface on the $Y$ axis from the data on the curved surface is derived.}
\end{figure}

Here, since $z_I$ and $z_J$ are on the cross-sectional curve, the following equation holds.

\begin{align}
z_I &= f(y_I) \notag\\
z_J &= f(y_J)
\tag{3-2}
\end{align}

The following equation is obtained from equation (2-15).

\begin{align}
&\Phi(k_x,\, y_I,\, y_J,\, k) \notag\\
&= \frac{1}{(2\pi)^2}
\int_{-\infty}^{\infty}\!\!\int_{-\infty}^{\infty}
e^{-i(k_{y_1} y_I + k_{y_2} y_J)}\,
a_{I,J}(k_x,\, k_{y_1},\, k_{y_2},\, k)\,
e^{is_1(k,\,k_x,\,k_{y_1},\,k_{y_2})\,z_I}\,
e^{is_2(k,\,k_x,\,k_{y_1},\,k_{y_2})\,z_J}\,
dk_{y_1}\,dk_{y_2}
\tag{3-3}
\end{align}

%%% ページ9 %%%

This formula can be interpreted and written as follows.

\begin{align}
&\Phi(k_x,\, y_I,\, y_J,\, k)\,\delta(y_1 - y_I)\,\delta(y_2 - y_J) \notag\\
&= \frac{1}{(2\pi)^2}
\int_{-\infty}^{\infty}\!\!\int_{-\infty}^{\infty}
e^{-i(k_{y_1} y_1 + k_{y_2} y_2)}\,
a_{I,J}(k_x,\, k_{y_1},\, k_{y_2},\, k)\,
e^{is_1(k,\,k_x,\,k_{y_1},\,k_{y_2})\,z_I}\,
e^{is_2(k,\,k_x,\,k_{y_1},\,k_{y_2})\,z_J}\,
dk_{y_1}\,dk_{y_2}
\tag{3-4}
\end{align}

The following equation is obtained by taking the Fourier transform of both sides.

\begin{align}
&\int_{-\infty}^{\infty}\!\!\int_{-\infty}^{\infty}
e^{i(k'_{y_1} y_1 + k'_{y_2} y_2)}\,
\Phi(k_x,\, y_I,\, y_J,\, k)\,\delta(y_1 - y_I)\,\delta(y_2 - y_J)\;dy_1\,dy_2
\notag\\
&= \int_{-\infty}^{\infty}\!\!\int_{-\infty}^{\infty}
e^{i(k'_{y_1} y_1 + k'_{y_2} y_2)}
\left\{
\frac{1}{(2\pi)^2}
\int_{-\infty}^{\infty}\!\!\int_{-\infty}^{\infty}
e^{-i(k_{y_1} y_1 + k_{y_2} y_2)}\,
a_{I,J}(k_x,\, k_{y_1},\, k_{y_2},\, k)\,
e^{is_1(k,\,k_x,\,k_{y_1},\,k_{y_2})\,z_I}\,
e^{is_2(k,\,k_x,\,k_{y_1},\,k_{y_2})\,z_J}\,
dk_{y_1}\,dk_{y_2}
\right\}dy_1\,dy_2
\tag{3-5}
\end{align}

%%% ページ10 %%%

When integrated,

\begin{align}
&e^{i(k'_{y_1} y_I + k'_{y_2} y_J)}\,\Phi(k_x,\, y_I,\, y_J,\, k) \notag\\
&= \left\{
\int_{-\infty}^{\infty}\!\!\int_{-\infty}^{\infty}
\delta(k_{y_1} - k'_{y_1})\,\delta(k_{y_2} - k'_{y_2})\,
a(k_x,\, k_{y_1},\, k_{y_2},\, k)\,
e^{is_1(k,\,k_x,\,k_{y_1},\,k_{y_2})\,z_I}\,
e^{is_2(k,\,k_x,\,k_{y_1},\,k_{y_2})\,z_J}\,
dk_{y_1}\,dk_{y_2}
\right\} \notag\\
&= a_{I,J}(k_x,\, k'_{y_1},\, k'_{y_2},\, k)\,
e^{is_1(k,\,k_x,\,k'_{y_1},\,k'_{y_2})\,z_I}\,
e^{is_2(k,\,k_x,\,k'_{y_1},\,k'_{y_2})\,z_J}
\tag{3-6}
\end{align}

In this way, $a_{I,J}$ are obtained as follows.

\begin{equation}
a_{I,J}(k_x,\, k_{y_1},\, k_{y_2},\, k)
= e^{i(k_{y_1} y_I + k_{y_2} y_J)}\,
e^{-is_1(k,\,k_x,\,k_{y_1},\,k_{y_2})\,z_I}\,
e^{-is_2(k,\,k_x,\,k_{y_1},\,k_{y_2})\,z_J}\,
\Phi(k_x,\, y_I,\, y_J,\, k)
\tag{3-7}
\end{equation}

The following equation is obtained by calculating the sum for all $I$ and $J$.

%%% ページ10 続き %%%

\begin{align}
a(k_x,\, k_{y_1},\, k_{y_2},\, k)
&= \sum_{I,J} a_{I,J}(k_x,\, k_{y_1},\, k_{y_2},\, k) \notag\\
&= \sum_{I,J}
e^{i(k_{y_1} y_I + k_{y_2} y_J)}\,
e^{-is_1(k,\,k_x,\,k_{y_1},\,k_{y_2})\,z_I}\,
e^{-is_2(k,\,k_x,\,k_{y_1},\,k_{y_2})\,z_J}\,
\Phi(k_x,\, y_I,\, y_J,\, k)
\tag{3-8}
\end{align}

In this way, it becomes possible to convert to the boundary condition at the plane $z = 0$.
Deriving the solution of the partial differential equation (2-7) with the boundary condition at
$z = 0$ gives the following from Eq.\ (2-15).

\begin{align}
&\phi(x,\, y_1,\, y_2,\, z_1,\, z_2,\, k) \notag\\
&= \frac{1}{(2\pi)^3}
\int_{-\infty}^{\infty}\!\!\int_{-\infty}^{\infty}\!\!\int_{-\infty}^{\infty}
e^{-i(k_x x + k_{y_1} y_1 + k_{y_2} y_2)}\,
a(k_x,\, k_{y_1},\, k_{y_2},\, k)\,
e^{is_1(k_x,\,k_{y_1},\,k_{y_2})\,z_1}\,
e^{is_2(k_x,\,k_{y_1},\,k_{y_2})\,z_2}\,
dk_x\,dk_{y_1}\,dk_{y_2}
\tag{3-9}
\end{align}

At this time, the reconstructed image can be obtained by integrating the following equation
with respect to $k$.

\begin{align}
&\phi(x,\, y,\, y,\, z,\, k)
= \lim_{y_1 \to y}\left[\phi(x,\, y_1,\, y,\, z,\, k)\right] \notag\\
&= \lim_{y_1 \to y}\left[
\frac{1}{(2\pi)^3}
\int_{-\infty}^{\infty}\!\!\int_{-\infty}^{\infty}\!\!\int_{-\infty}^{\infty}
e^{-i(k_x x + k_{y_1} y_1 + k_y y)}\,
a(k_x,\, k_{y_1},\, k_y,\, k)\,
e^{i\sqrt{\left(\sqrt{k^2 - k_{y_1}^2} + \sqrt{k^2 - k_y^2}\right)^2 - k_x^2}\;z}\,
dk_x\,dk_{y_1}\,dk_y
\right]
\tag{3-10}
\end{align}

A variable $k_z$ is introduced as follows. The formula for expressing $k$ by $k_z$ and the function
obtained by differentiating it are also shown.

\begin{align}
k_z &= \sqrt{\left(\sqrt{k^2 - k_{y_1}^2} + \sqrt{k^2 - k_{y_2}^2}\right)^2 - k_x^2}
\notag\\[6pt]
k &= \frac{1}{2}\sqrt{k_x^2 + k_z^2 + \frac{(k_{y_1}^2 - k_{y_2}^2)^2}{k_x^2 + k_z^2} + 2(k_{y_1}^2 + k_{y_2}^2)}
\notag\\[6pt]
\frac{dk}{dk_z} &= \frac{k_z\,\sqrt{k^2 - k_{y_1}^2}\,\sqrt{k^2 - k_{y_2}^2}}{k\,(k_x^2 + k_z^2)}
\tag{3-11}
\end{align}

Finally, the image function $\rho(x,\, y,\, z)$ of the reconstruction result is described as follows.

%%% ページ11 %%%

\begin{align}
\rho(x,\, y,\, z)
&= \int_{-\infty}^{\infty} \phi(x,\, y,\, y,\, z,\, k)\,dk
\notag\\[6pt]
&= \int_{-\infty}^{\infty}
\lim_{y_1 \to y}\left[\phi(x,\, y_1,\, y,\, z,\, k)\right]dk
\notag\\[6pt]
&= \int_{-\infty}^{\infty}
\lim_{y_1 \to y}\left[
\frac{1}{(2\pi)^3}
\int_{-\infty}^{\infty}\!\!\int_{-\infty}^{\infty}\!\!\int_{-\infty}^{\infty}
e^{-i(k_x x + k_{y_1} y_1 + k_y y)}\,
a(k_x,\, k_{y_1},\, k_y,\, k)\,
e^{i\sqrt{\left(\sqrt{k^2 - k_{y_1}^2} + \sqrt{k^2 - k_y^2}\right)^2 - k_x^2}\;z}\,
dk_x\,dk_{y_1}\,dk_y
\right]dk
\notag\\[6pt]
&= \int_{-\infty}^{\infty}
\lim_{y_1 \to y}\left[
\frac{1}{(2\pi)^3}
\int_{-\infty}^{\infty}\!\!\int_{-\infty}^{\infty}\!\!\int_{-\infty}^{\infty}
e^{-i(k_x x + k_{y_1} y_1 + k_y y)}\,
e^{ik_z z}\,
a(k_x,\, k_{y_1},\, k_y,\, k)
\left(\frac{dk}{dk_z}\right)
dk_x\,dk_{y_1}\,dk_y\,dk_z
\right]
\tag{3-12}
\end{align}

The theory for the inverse problem of wave scattering in this paper has a wide range of practical
implications such as medical imaging and various non-destructive subsurface imaging.
Although not mentioned in this paper about\quad the details of experimental setups, the theory that
clearly incorporates the physical characteristics of the medium and the structure of the antenna
actually used is equipped into the device developed by our research group. As the future work,
important improvements should be made according to various application targets based on this
theory.

\vspace{2em}

%%% ページ12（謝辞）%%%

This work was supported by the Research Institute for Mathematical Sciences, an International
Joint Usage/Research Center located in Kyoto University, and Medical Research and
Development Programs Focused on Technology Transfers: Development of Advanced
Measurement and Analysis Systems (AMED-SENTAN) in Japan Agency for Medical Research
and Development, and Research on Development of New Medical Devices in Japan Agency
for Medical Research and Development, Research on Development of New Medical Devices
in Japan Agency for Medical Research and Development, and Grant-in-Aid for Transformative
Research Areas in Ministry of Education, Culture, Sports, Science and Technology -- Japan.

\end{document}